\documentclass{beamer}
\usepackage{ctex}
\usetheme{SimplePlus}
\setbeamertemplate{footline}{}

\usepackage{amsmath,amssymb,mathtools}
\usepackage{tikz}
\usetikzlibrary{arrows.meta}
\usepackage{tikz-cd}

\usepackage{unicode-math}
%\setmathfont{STIX Two Math}
%\setmathfont{TeX Gyre Termes Math}
%\setmathfont{Libertinus Math}
%\setmathfont{Fira Math}

\usepackage[
  hybrid,                            % 允许在 Markdown 中混用 LaTeX 命令
  texMathDollars                     % 启用 $...$ 和 $$...$$ 数学公式
]{markdown}

\usepackage{graphicx} % Allows including images
\graphicspath{{./images/}}

\title{End}

\date{}

\begin{document}

\begin{frame}
    % Print the title page as the first slide
    \titlepage
\end{frame}

\begin{frame}[fragile,shrink]{dinatural transformation}
Let $F, G \colon C^{\mathrm{op}} \times C \to D$ be functors. A \textbf{dinatural transformation} from $F$ to $G$, sometimes written
\[
\alpha \colon F \overset{\bullet}{\to} G,
\]
consists of a collection of morphisms
\[
\alpha_c \colon F(c,c) \to G(c,c)
\]
such that for every morphism $f \colon c \to c'$ in $C$,
\[
G(c,f)\alpha_c F(f,c) = G(f,c')\alpha_{c'} F(c',f) \colon F(c',c) \to G(c,c')
\]
If drawn out as a commutative diagram, this becomes a ``hexagon identity'':

\begin{center}
\begin{tikzcd}[row sep=2.2em, column sep=2.2em]
  & F(c,c) \arrow[r,"{\alpha_c}"] & G(c,c) \arrow[dr,"{G(c,f)}"] & \\
  F(c',c) \arrow[ur,"{F(f,c)}"] \arrow[dr,"{F(c',f)}"'] & & & G(c,c') \\
  & F(c',c') \arrow[r,"{\alpha_{c'}}"'] & G(c',c') \arrow[ur,"{G(f,c')}"'] &
\end{tikzcd}
\end{center}
\end{frame}

\begin{frame}{Special cases}
If $F$ and $G$ both factor through the projection $C^{\mathrm{op}}\times C\to C$, then the notion reduces to an ordinary natural transformation, and similarly if they both factor through $C^{\mathrm{op}}$.

If $F$ factors through $C$ while $G$ factors through $C^{\mathrm{op}}$, then we obtain a notion of natural transformation from a covariant functor to a contravariant one, and dually.

If $G$ is constant, the hexagon identity reduces to the ``domain'' version of extranaturality, involving a commutative square of the form
\[
\alpha_c F(f,c)=\alpha_{c'}F(c',f)\colon F(c',c)\to G
\]
where $G$ is constant with respect to the argument $c$. Similarly, if $F$ is constant, it yields the ``codomain'' version with a commutative square of the form
\[
G(c,f)\alpha_c=G(f,c')\alpha_{c'}\colon F\to G(c,c')
\]
when $F$ is constant with respect to the argument $c$.
\end{frame}

\begin{frame}[fragile,shrink]{Wedge}
设：
$$
H:\mathcal C^{op}\times\mathcal C\to\mathcal E.
$$

以及 constant bifunctor：
$$
\Delta X:\mathcal C^{op}\times\mathcal C\to\mathcal E
$$

定义：
$$
\Delta X(c,c')=X.
$$

那么：
$$
\boxed{
\omega:\Delta X\Rrightarrow H
}
$$

就是一个 dinatural transformation。

它的 component 是：
$$
\omega_c:X\to H(c,c).
$$

dinaturality condition 变成：
$$
\boxed{
H(f,1_{c'})\circ\omega_{c'}
=
H(1_c,f)\circ\omega_c
}
$$

这正是 \textbf{wedge} 的定义。

所以：
$$
\boxed{
\text{wedge}
=
\text{dinatural transformation }
\Delta X\Rrightarrow H
}
$$
\end{frame}

\begin{frame}{Cowedge}
如果：
$$
\iota:H\Rrightarrow\Delta X,
$$

那么 component：
$$
\iota_c:H(c,c)\to X.
$$

dinaturality condition 变成：
$$
\boxed{
\iota_c\circ H(1_c,f)
=
\iota_{c'}\circ H(f,1_{c'})
}
$$

这就是 \textbf{cowedge}。

所以：
$$
\boxed{
\text{cowedge}
=
\text{dinatural transformation }
H\Rrightarrow\Delta X
}
$$
\end{frame}

\begin{frame}[shrink]{end / coend}
给定：
$$
H:\mathcal C^{op}\times\mathcal C\to\mathcal E.
$$

\textbf{End}

一个 end 是一个 universal dinatural transformation：
$$
\boxed{
\omega:
\Delta\left(\int_cH(c,c)\right)
\Rrightarrow H
}
$$

即：
$$
\omega_c:
\int_cH(c,c)\to H(c,c).
$$

它 universal 的意思是：

对于任意：
$$
\alpha:\Delta X\Rrightarrow H,
$$

存在唯一：
$$
u:X\to\int_cH(c,c)
$$

使：
$$
\boxed{
\alpha=\omega\circ\Delta u
}
$$
\end{frame}

\begin{frame}
\textbf{Coend}

一个 coend 是一个 universal dinatural transformation：
$$
\boxed{
\iota:
H\Rrightarrow
\Delta\left(\int^cH(c,c)\right)
}
$$

即：
$$
\iota_c:
H(c,c)\to\int^cH(c,c)
$$

对于任意：
$$
\alpha:H\Rrightarrow\Delta X,
$$

存在唯一：
$$
u:\int^cH(c,c)\to X
$$

使：
$$
\boxed{
\alpha=\Delta u\circ\iota
}
$$
\end{frame}

\begin{frame}{典型例子}
\begin{itemize}
  \item end 给出自然变换
  \item Yoneda 引理的 end 形式
  \item Yoneda 引理的 coend 形式
  \item Kan extension
  \item Geometric realization
\end{itemize}
\end{frame}

\begin{frame}{优点}
\begin{itemize}
  \item End/coend 最大的优势：把“自然性条件”打包掉。
  \item End/coend 将 universal construction 转化成一种可进行代数操作的表达式。
\end{itemize}

$$
\boxed{
\textbf{End/coend 是范畴论的“积分/求和演算”。}
}
$$

而 \textbf{Kan extension 更像是你要解决的 universal problem}：
$$
\boxed{
\text{Kan extension：我要构造什么？}
}
$$
$$
\boxed{
\text{End/coend：这个 universal object 如何把它写成一个可操作的公式？}
}
$$
\end{frame}

% 参考文献帧，长列表自动分页
\begin{frame}[allowframebreaks]{参考文献}
\begin{thebibliography}{99}

\bibitem{ncatlab_dinatural}
\url{https://ncatlab.org/nlab/show/dinatural+transformation}

\bibitem{end_natural_transformation}
\url{https://chatgpt.com/s/t_6a96dbb328ec81919d35598bf6fee18d}

\end{thebibliography}
\end{frame}
\end{document}